\documentclass[11pt,a4paper]{article}
\usepackage[utf8]{inputenc}
\usepackage[T1]{fontenc}
\usepackage[brazil]{babel}
\usepackage{lmodern,amsmath,amssymb,amsthm,mathtools}
\usepackage[margin=2.25cm]{geometry}
\usepackage{booktabs,tabularx,array}
\usepackage{microtype,xcolor,enumitem,xurl,hyperref}
\hypersetup{colorlinks=true,linkcolor=blue!45!black,citecolor=blue!45!black,urlcolor=blue!55!black,
 pdftitle={Por que o treinamento de HaxBall passou a funcionar melhor}}
\setlength{\emergencystretch}{3em}
\setlist{itemsep=2pt,topsep=4pt}
\newcommand{\E}{\mathbb{E}}
\newcommand{\R}{\mathbb{R}}
\newcommand{\ind}{\mathbf{1}}
\newcommand{\norm}[1]{\left\lVert#1\right\rVert}
\newcommand{\code}[1]{\nolinkurl{#1}}
\newcolumntype{Y}{>{\raggedright\arraybackslash}X}
\newtheorem{proposicao}{Proposição}
\newtheorem{observacao}{Observação}
\title{\textbf{Por que o treinamento de HaxBall\\passou a funcionar melhor?}\\[5pt]
\large Do HaxBall a outros jogos e à robótica: uma leitura guiada sobre aprender com consequências}
\author{Relatório técnico do projeto \texttt{haxball-rl}}
\date{15 de outubro de 2026}
\begin{document}
\maketitle
\begin{abstract}
Este artigo explica por que versões anteriores de um agente de HaxBall apresentavam progresso frágil e o que mudou na configuração que produziu melhora mensurável. A tese central é que não bastava aumentar a duração do treinamento: era necessário tornar coerentes o incentivo, a estimativa das consequências, os adversários e o critério de avaliação. A partir de código e registros do projeto, apresentam-se a formulação de aprendizado por reforço, a identidade matemática de recompensas por potenciais, o tratamento de cortes artificiais, a interação entre ator e crítico e a não estacionariedade do self-play. Cada mecanismo é acompanhado de uma interpretação para leitores sem formação em aprendizado de máquina. Um benchmark reservado do checkpoint 2.550 mostra score de 0,66895 contra sua origem em estados aleatórios com ações amostradas; a continuação GPU também apresenta avanço descritivo, embora oscilatório. Gráficos distinguem resultados observados de exemplos teóricos. As intervenções foram combinadas e não permitem atribuição causal individual. ``Dar certo'' significa aqui obter ganho verificável nas condições testadas, sem afirmar solução definitiva do jogo. O teorema de Ng, Harada e Russell é explicado por exemplos de percursos e contabilidade de recompensas. Aplicações a Pokémon, robótica e outros jogos são propostas conceituais, sem experimentos realizados nesses domínios. A cooperação 2v2 permanece uma perspectiva ainda não validada.
\end{abstract}
\noindent\textbf{Palavras-chave:} aprendizado por reforço; PPO; reward shaping; HaxBall; cooperação multiagente.

\section{Introdução e escopo da evidência}
A melhora de um agente não deve ser inferida apenas da recompensa acumulada. É necessário verificar se ele vence adversários relevantes, preserva competências e generaliza para situações não usadas na seleção. Em linguagem simples: aprender a ganhar pontos durante o treino e aprender a jogar bem são avaliações relacionadas, mas distintas.

Este estudo combina inspeção de código, métricas do treino até $10\,000$ iterações, benchmarks anteriores e um teste de reset 2v2. A referência de código é a branch \code{main}, revisão \code{9ea2497}; o diagnóstico usa o estado do projeto em \code{c48c249} e os dados GPU publicados em \code{2694aa8}. O checkpoint selecionado é o da iteração $9\,900$.

A comparação com a \code{main} descreve mudanças de engenharia. Não houve novo experimento com ambas as branches treinadas do zero, sob o mesmo orçamento. O agente denominado \emph{origem} nos benchmarks anteriores é o início da nossa continuação, não uma reprodução independente da \code{main}. Consequentemente, ganhos não podem ser atribuídos isoladamente a uma alteração.


\paragraph{Como ler o artigo.}
As fórmulas são acompanhadas por interpretação verbal. Um \emph{checkpoint} é uma versão salva do agente; uma \emph{iteração} reúne coleta e atualização; uma \emph{seed} fixa uma sequência pseudoaleatória. Seeds de avaliação variam as provas; seeds de treino repetem o processo de aprendizado. \emph{Greedy} escolhe sempre a ação mais provável; \emph{amostrado} sorteia de acordo com as probabilidades aprendidas. O score usado adiante vale um por vitória, meio por empate e zero por derrota: não é a própria recompensa de treinamento.


\subsection*{Uma leitura em três voltas}
Este texto pode ser lido sem tentar dominar todas as equações de uma vez. Na primeira volta, acompanhe as histórias e figuras: que comportamento queremos, que orientação oferecemos e como verificamos o resultado? Na segunda, observe as três grandezas centrais: recompensa é o que acabou de acontecer; valor é a previsão do que ainda virá; vantagem é o quanto a decisão foi melhor ou pior do que se esperava. Na terceira, volte às equações para conferir por que as afirmações fazem sentido.

O símbolo $\Phi$, que aparecerá muitas vezes, pode ser lido simplesmente como ``nota da situação''. O símbolo $\gamma$ diz quanto o futuro pesa. A palavra \emph{política} significa a maneira de agir do agente, não uma regra que escrevemos manualmente para ele seguir. Não é necessário memorizar os hiperparâmetros para compreender o argumento.

\section{O que significa dizer que antes deu errado?}
\subsection{O problema era a fragilidade, não a ausência total de aprendizado}
A expressão ``antes deu errado'' exige precisão. As versões anteriores executavam atualizações válidas e chegaram a aprender comportamentos úteis. O problema documentado era que melhorar uma métrica não assegurava preservar outras competências. Houve oscilação entre checkpoints, diferença entre jogo amostrado e jogo determinístico e perda contra uma referência neural apesar de ganhos contra heurísticas.

Na continuação noturna, o melhor painel interno marcou $0{,}87147$, enquanto o endpoint caiu para $0{,}70907$. Em uma avaliação posterior, com condições próprias, o score greedy contra a referência neural passou de $0{,}90854$ na origem para $0{,}58959$ no melhor candidato. Esses dois exemplos vêm de painéis distintos e não devem ser unidos numa única curva. Eles demonstram dois problemas diferentes: regressão ao continuar treinando e ganho agregado que encobre perda contra um adversário.

Também houve uma tentativa de corrigir apenas o bootstrap de timeout, comparando três pares de seeds B0/E1. A diferença média E1 menos B0 foi $+0{,}01238$ em estados aleatórios, mas $-0{,}05368$ em kickoff. Uma correção compatível com um objetivo matemático não garantiu melhorar o score de partida. Isso foi informação útil: não existia evidência para continuar alterando a mesma máscara até aparecer um número favorável.

\subsection{Uma cadeia de aprendizado, com vários pontos de falha}
A política é uma regra probabilística de decisão; o crítico é uma previsão de consequências. O ambiente fornece experiências e a avaliação verifica o resultado. A Figura~\ref{fig:chain} organiza essas funções.
\begin{figure}[htbp]\centering\small
\begin{tabular}{c}
\fbox{\parbox{.86\linewidth}{\centering\textbf{Ambiente e adversário}\\Determinam quais situações e consequências o agente encontra.}}\\[3pt]
$\downarrow$\\[3pt]
\fbox{\parbox{.86\linewidth}{\centering\textbf{Recompensa e crítico}\\Definem o incentivo e estimam o mérito futuro das decisões.}}\\[3pt]
$\downarrow$\\[3pt]
\fbox{\parbox{.86\linewidth}{\centering\textbf{Atualização do ator}\\Aumenta ou reduz probabilidades das ações, com passos limitados.}}\\[3pt]
$\downarrow$\\[3pt]
\fbox{\parbox{.86\linewidth}{\centering\textbf{Avaliação externa à atualização}\\Testa adversários, lados e modos; detecta regressões e escolhe o checkpoint.}}
\end{tabular}
\caption{Esquema conceitual, não um resultado experimental. As experiências do próximo ciclo dependem da política atual; uma falha em qualquer ligação pode comprometer o conjunto. Avaliar corretamente melhora a seleção e o diagnóstico, mas não fornece por si só um gradiente melhor.}\label{fig:chain}
\end{figure}

Uma analogia é treinar um jogador com exercícios, um placar e um treinador. Se o placar recompensa a tarefa errada, ele aprende o exercício errado. Se a previsão do treinador é ruim, a orientação pode ser ruim mesmo com um placar correto. Se só enfrenta um estilo, pode se especializar. Se a prova só mede esse estilo, a especialização parece uma melhora universal. A intervenção estrutural tratou conjuntamente esses pontos.

\section{Fundamentos de RL: como uma experiência vira aprendizado}
\subsection{Retorno, valor e vantagem em linguagem operacional}
O aprendizado por reforço não recebe para cada imagem uma ação humana correta. Recebe consequências depois de agir. Para desconto $\gamma$, o retorno é
\begin{equation}
 G_t=r_t+\gamma r_{t+1}+\gamma^2r_{t+2}+\cdots.
\end{equation}
O fator $\gamma$ reduz o peso do futuro. O valor $V^\pi(s)=\E[G_t\mid s_t=s]$ resume o futuro médio ao seguir a política $\pi$. Já $Q^\pi(s,a)$ considera primeiro uma ação específica. A vantagem ideal é
\begin{equation}
 A^\pi(s,a)=Q^\pi(s,a)-V^\pi(s).
\end{equation}
Ela compara uma decisão com o que seria esperado naquela situação. Uma ação com vantagem positiva não precisa resultar imediatamente em gol: pode apenas criar uma oportunidade melhor do que a alternativa habitual. Uma vantagem negativa também não significa necessariamente derrota: pode representar uma oportunidade desperdiçada.

Sob as hipóteses usuais do gradiente de política, sua direção tem a forma
\begin{equation}
 \nabla_\theta J\ \propto\
 \E\!\left[\nabla_\theta\log\pi_\theta(a_t\mid o_t)\,A^\pi(s_t,a_t)\right].
 \label{eq:pg}
\end{equation}
O fator de proporcionalidade depende da convenção de normalização da distribuição descontada de estados. A leitura prática é: reforçar decisões estimadas como melhores do que o esperado e enfraquecer as piores. Portanto, se a estimativa de vantagem estiver distorcida, atualizar corretamente a rede ainda pode reforçar ações inadequadas.

Um baseline $b(s)$ independente da ação não muda a esperança do gradiente ideal, pois
\begin{equation}
 \E_{a\sim\pi_\theta}[\nabla_\theta\log\pi_\theta(a\mid s)b(s)]
 =b(s)\nabla_\theta\sum_a\pi_\theta(a\mid s)=0.
\end{equation}
Isso explica por que subtrair o valor pode reduzir variância sem alterar o objetivo ideal. Não significa que um crítico aproximado seja inofensivo: o GAE usa bootstrap, há amostras finitas, e PPO normaliza e reutiliza vantagens em uma otimização aproximada.

\subsection{Por que aumentar o número de partidas não resolve tudo}
Mais dados reduzem algumas fontes de ruído e permitem mais atualizações. Contudo, repetir experiências de uma distribuição limitada não cria automaticamente novos desafios. Uma recompensa inadequada também não se torna adequada após milhões de repetições. Há três perguntas diferentes: o algoritmo está executando corretamente? Está otimizando o objetivo definido? Esse objetivo e essa distribuição produzem a habilidade desejada? Os testes de software respondem principalmente à primeira; partidas reservadas ajudam a responder à terceira.

\section{Modelo do problema e política}
Sejam $s_t$ o estado físico, $o_t^i=O_i(s_t)$ a observação do jogador $i$ e $\pi_j^-$ uma política adversária congelada. Para uma distribuição curricular $q(j)$, o objetivo de referência é
\begin{equation}
 J(\theta;q)=\E_{j\sim q,\,\pi_\theta,\,\pi_j^-}
 \left[\sum_{t=0}^{\infty}\gamma^t r_t^{\mathrm{gol}}\right],
 \qquad r_t^{\mathrm{gol}}=25\bigl(g_t^+-g_t^-\bigr),
 \label{eq:objective}
\end{equation}
com $g_t^+,g_t^-\in\{0,1\}$ indicando gol a favor e contra. A observação atual é vetorial, sem memória explícita. Posições e velocidades não incluem necessariamente todo o estado interno relevante, como o rearme do chute; não se assume observabilidade Markov plena.

As ações e sua distribuição são
\begin{equation}
 a_t^i=(a_x,a_y,a_k)\in\{0,1,2\}^2\times\{0,1\},\qquad
 \pi_\theta(a\mid o)=\prod_{h\in\{x,y,k\}}\pi_{\theta,h}(a_h\mid o).
\end{equation}
Os bins de movimento representam $-1,0,+1$; o chute representa soltar ou pressionar. O wrapper de execução rearma o chute dentro da janela física. Pressionar não comprova um chute efetivo.

O ator possui duas camadas de $256$ unidades com $\tanh$; o crítico atual tem parâmetros separados. Com decisões a cada oito ticks e motor a $60$ Hz,
\begin{equation}
 \Delta t=\frac{8}{60}\,\mathrm{s},\qquad
 \gamma=2^{-\Delta t/12}\approx0{,}99232795.
\end{equation}
Para $N=512$ ambientes, $T=64$ decisões e $m$ jogadores aprendizes por ambiente,
\begin{equation}
 B_{\mathrm{agente}}=NTm,\qquad B_{\mathrm{mundo}}=NT,
 \qquad T\Delta t\approx8{,}53\,\mathrm{s}.
 \label{eq:budget}
\end{equation}
No 1v1, $m=1$ e são $32\,768$ transições do aprendiz por iteração. Contar dois companheiros em 2v2 dobra transições de agentes, mas não cria o dobro de partidas independentes.

\section{Recompensa: formulação e condições}
Sejam $b$ a posição da bola, $p_i$ a posição do jogador, $c_i^+$ o centro do gol atacado e $\mathcal G_i^\pm$ os segmentos de gol. Definem-se
\begin{align}
 u_i(s)&=\frac{c_i^+-b}{\max(\norm{c_i^+-b},\varepsilon)},&
 q_i(s)&=b-22u_i(s),\\
 \Phi_{\mathrm{pos}}^i(s)&=-\frac{\operatorname{clip}(\norm{p_i-q_i(s)}-2,\,0,\,2000)}{\ell},&
 \Phi_{\mathrm{bola}}^i(s)&=\frac{d(b,\mathcal G_i^-)-d(b,\mathcal G_i^+)}{2\ell},
\end{align}
onde $\ell=21{,}12$ é a escala de deslocamento por decisão no mapa de referência. Para um segmento $[v_0,v_1]$,
\begin{equation}
 d(b,[v_0,v_1])=\norm{b-v_0-\operatorname{clip}\!\left(
 \frac{(b-v_0)^\top(v_1-v_0)}{\norm{v_1-v_0}^2},0,1\right)(v_1-v_0)}.
\end{equation}
O potencial combinado e a recompensa efetiva são
\begin{align}
 \Phi^i(s)&=0{,}15\Phi_{\mathrm{pos}}^i(s)+0{,}08\Phi_{\mathrm{bola}}^i(s),\\
 r_t^i&=25(g_t^+-g_t^-)+\gamma\bar\Phi^i(s_{t+1})-\Phi^i(s_t),
 \label{eq:shaping}
\end{align}
com $\bar\Phi=0$ em terminal de gol. Os pesos numéricos foram preservados, mas as funções auxiliares mudaram: posicionamento atrás da bola substituiu o deslocamento simples em sua direção, e o progresso passou a comparar os dois gols.


\subsection{Por que a estrutura da recompensa importa mais que o nome dos pesos}
Dizer que os pesos continuam $0{,}15$, $0{,}08$ e $25$ não significa que a tarefa de aprendizado continuou igual. Um peso multiplica uma função; mudar a função muda o sinal de treinamento. Na tabela corrected intermediária, perseguir a bola era um incentivo mais direto; essa descrição não reconstrói por si só toda a recompensa padrão da dependência usada pela main histórica. Agora, o potencial de posição avalia um ponto atrás dela, alinhado ao gol. Aproximar-se da bola pela frente e aproximar-se por uma posição de chute deixam de ser equivalentes para o sinal auxiliar.

Uma recompensa densa arbitrária pode modificar a preferência entre estratégias. Para uma rota $\mathcal A$ e outra $\mathcal B$, pode ocorrer
\begin{equation}
 G_{\rm gol}(\mathcal A)>G_{\rm gol}(\mathcal B),\quad
 G_{\rm gol}(\mathcal A)+G_{\rm aux}(\mathcal A)
 <G_{\rm gol}(\mathcal B)+G_{\rm aux}(\mathcal B).
\end{equation}
Nesse caso, o agente que maximiza a soma pode preferir a pior rota para o objetivo de gols. Isso é um risco geral de desenho de incentivos, não uma demonstração de que o bot antigo explorou exatamente esse mecanismo. O PBRS busca oferecer orientação local com um contrato matemático para a soma completa.

\subsection{O teorema de Ng, Harada e Russell: orientar sem trocar o objetivo}
O resultado de Andrew Ng, Daishi Harada e Stuart Russell, publicado em 1999, responde a uma pergunta prática: é possível oferecer pistas intermediárias sem mudar quais políticas são ótimas? Em um processo de decisão de Markov (MDP), o estado reúne a informação necessária para prever o próximo passo, dadas as ações. Para esse modelo com desconto, uma construção que preserva essas políticas é acrescentar $F(s,s')=\gamma\Phi(s')-\Phi(s)$ à recompensa original, com as condições de horizonte e fronteira adequadas \cite{ng1999}. A conta abaixo mostra a ideia central. Não é uma afirmação de que qualquer bônus intuitivamente razoável tenha essa propriedade.

\paragraph{A conta que explica o resultado.}

Para potencial limitado, desconto $0<\gamma<1$ e fronteiras consistentes,
\begin{equation}
 \sum_{t=0}^{H-1}\gamma^t[\gamma\Phi(s_{t+1})-\Phi(s_t)]
 =-\Phi(s_0)+\gamma^H\Phi(s_H).
\end{equation}
Em terminal absorvente com potencial zero, ou no limite $H\to\infty$, resta um deslocamento independente das ações, condicionado a $s_0$.
A expressão parece abstrata, mas descreve um cancelamento simples. O potencial do estado seguinte entra uma vez como crédito e, na decisão seguinte, volta como débito. Os fatores de desconto fazem esses valores se alinharem. Sobram apenas a situação inicial e a fronteira final. Quando a fronteira final vale zero, a soma auxiliar depende só de onde a trajetória começou.
Essa identidade fundamenta o PBRS \cite{ng1999}. A referência de invariância é o objetivo esparso de gols, não a tabela antiga. Não se segue que PPO com função aproximadora convirja ou aprenda mais rapidamente.


\paragraph{Um exemplo numérico da demonstração.}
Considere, apenas para ilustração, $\gamma=0{,}9$ e potenciais $(-10,-6,-2,0)$ ao longo de três decisões. Os sinais auxiliares são $4{,}6$, $4{,}2$ e $2$. O retorno auxiliar descontado é
\begin{equation}
 4{,}6+0{,}9\times4{,}2+0{,}9^2\times2=10=-\Phi(s_0).
\end{equation}
Outra trajetória que comece com potencial $-10$ e termine no mesmo terminal absorvente também produz retorno auxiliar $10$, com as condições da proposição. Assim, o potencial redistribui a orientação no tempo sem acrescentar preferência líquida entre essas trajetórias; os gols continuam decidindo seu mérito no objetivo de referência. Os números deste exemplo não são dados do treinamento.

Como uma redistribuição sem vantagem líquida pode ajudar? Porque o crítico e o ator aprendem com dados finitos, aproximações e passos locais. Receber informação em estados intermediários pode facilitar a estimativa do mérito, embora não haja garantia. O objetivo ideal pode ser equivalente enquanto a dificuldade prática da otimização muda.


\subsubsection*{A intuição: um adiantamento que precisa fechar a conta}
Imagine que o objetivo seja entregar uma encomenda. Receber dinheiro toda vez que se passa por uma rua perto do destino pode incentivar dar voltas nessa rua. Já um adiantamento vinculado à posição atual precisa ser acertado quando a posição muda. Avançar libera orientação; desfazer o avanço retira essa orientação. O agente não deveria aumentar o retorno total apenas oscilando.

Essa analogia tem um detalhe essencial: o dinheiro do futuro vale menos por causa de $\gamma$. Por isso a fórmula não é simplesmente ``nota nova menos nota antiga''. O desconto precisa entrar no acerto. A Figura~\ref{fig:ngcycle} mostra uma conta pequena que pode ser conferida à mão.

\begin{figure}[htbp]\centering
\fbox{\parbox{.88\linewidth}{
\textbf{Exemplo didático: sair e voltar ao mesmo lugar.}\\[4pt]
$A\ (\Phi=0)\quad\longrightarrow\quad B\ (\Phi=1)\quad\longrightarrow\quad A\ (\Phi=0)$\\[5pt]
Com $\gamma=0{,}9$, usar apenas a diferença sem desconto dá sinais $+1$ e $-1$. A soma descontada é $1+0{,}9(-1)=0{,}1$. A volta aparentemente neutra produziu lucro auxiliar.\\[5pt]
Com PBRS, os sinais são $+0{,}9$ e $-1$. A soma descontada é $0{,}9+0{,}9(-1)=0$. O crédito e a devolução se compensam.
}}
\caption{Exemplo inventado para explicar o papel do desconto, não uma trajetória observada do bot. O estado inicial e o final têm potencial zero. Se não tivessem, o termo de fronteira também precisaria entrar na conta.}\label{fig:ngcycle}
\end{figure}

O exemplo não diz que toda repetição seja inútil. Um robô pode precisar recuar para contornar um obstáculo, e um jogador pode precisar devolver a bola para criar espaço. A questão é que o incentivo auxiliar não deveria transformar um ciclo sem benefício para a tarefa em uma fonte autônoma de retorno.

\subsubsection*{Por que a melhor decisão permanece a mesma?}
Chamemos de $G'$ o retorno com shaping e de $G$ o retorno original. Com as condições já explicadas,
\begin{equation}
 G'=G-\Phi(s_0),\qquad
 Q'^{\pi}(s,a)=Q^{\pi}(s,a)-\Phi(s).
\end{equation}
Para quem decide no mesmo estado $s$, a quantidade subtraída é igual para todas as ações. Subtrair dez pontos de todas as alternativas não muda qual delas tem a maior nota. Logo,
\begin{equation}
 \arg\max_a Q'^*(s,a)=\arg\max_a Q^*(s,a).
\end{equation}
É isso que significa preservar a política ótima nesse contexto. Não significa que as recompensas de cada passo fiquem iguais, nem que os valores numéricos de $Q$ sejam os mesmos. Significa preservar a preferência ideal entre maneiras de agir.

Essa conclusão pressupõe que as alternativas sejam comparadas a partir da mesma situação inicial, ou de uma distribuição inicial fixa. Se começamos a selecionar estados iniciais de forma diferente entre tratamentos, já modificamos outra parte do experimento. Também é necessário usar o mesmo desconto na recompensa auxiliar e no objetivo.

\paragraph{Isso significa que só existe uma transformação permitida?}
Não. Multiplicar todo o retorno por uma constante positiva também preserva a ordem das alternativas. A relevância do resultado é fornecer uma estrutura geral de shaping que não depende de adivinhar toda a dinâmica. Em um ambiente específico, outras mudanças podem por acaso preservar a solução; isso exige outra justificativa. Em particular, somar um bônus por passo não é sempre neutro: se os episódios têm durações diferentes, pode passar a valer a pena prolongá-los. Não se deve confundir essa operação com subtrair a mesma constante de todas as alternativas em um estado.

\subsubsection*{E se a nossa nota da situação estiver errada?}
O potencial não precisa ser um oráculo. Um mapa de distância que ignora obstáculos pode julgar boa uma posição que, na prática, é uma armadilha. Ainda assim, um potencial fixo, limitado e aplicado com as fronteiras corretas mantém a identidade de retorno. O problema é prático: as pistas podem conduzir a exploração para lugares ruins e tornar o aprendizado lento. Preservar o destino matemático não garante fornecer um bom caminho para chegar até ele.

No HaxBall, ficar atrás da bola parece uma pista útil para criar uma situação de chute. Não significa que ficar atrás seja sempre a melhor decisão tática. Às vezes a boa ação é defender, interceptar ou se afastar. O objetivo de gols é que deve decidir a estratégia final; o potencial tenta facilitar a orientação local.

Também se pode construir um potencial a partir de um modelo que prevê chance de sucesso. Se ele for congelado e satisfizer as condições, continua sendo um potencial. Se o atualizarmos durante o treino, a justificativa de potencial fixo já não pode ser usada automaticamente: seria necessário analisar a dependência temporal e os respectivos termos. Uma ideia promissora ainda precisa de contrato de implementação.

\subsubsection*{O que esse resultado não promete}
O teorema trata da relação entre objetivos, não do sucesso garantido de qualquer programa que use sua fórmula. Nossa rede pode representar mal uma estratégia, explorar pouco, estimar valores incorretamente ou gastar todo o orçamento antes de encontrar uma solução. O clipping do PPO e a coleta finita também fazem parte da dificuldade prática.

Um potencial baseado em observação incompleta pode continuar obedecendo à conta telescópica se for aplicado consistentemente; isso não devolve à política a informação que ela nunca recebeu. A identidade algébrica não exige que a nota seja uma descrição completa do mundo. Já a interpretação clássica de política ótima em estados usa o modelo de decisão especificado. Memória, informação oculta e restrições da rede continuam relevantes.

O teorema também não escolhe para nós a recompensa original. Se ela expressa mal o que desejamos, preservar suas soluções ótimas preserva o problema. No HaxBall, retorno descontado de gols e score de uma partida finita são critérios diferentes; PBRS não elimina essa diferença. Em um episódio realmente finito, não se pode simplesmente descartar um potencial final dependente das ações e continuar alegando invariância.

\paragraph{Uma pergunta para levar à próxima leitura.}
Quando surgir a ideia de ``dar um pontinho por uma coisa boa'', vale perguntar: esse ponto representa o objetivo que quero mudar, uma pista cuja conta fecha, ou apenas uma métrica que desejo observar? São três usos legítimos, mas com consequências diferentes.

\begin{observacao}
Para $\Phi(s)<0$, um estado estacionário pode gerar $(\gamma-1)\Phi(s)>0$. O sinal de um passo isolado não demonstra exploração indevida da recompensa; importam a soma descontada e suas fronteiras.
\end{observacao}

\subsection{Gols, timeouts e vantagens}
Sejam $z_t$ um término real e $\tau_t$ um corte artificial. O código usa a observação anterior ao autoreset para calcular
\begin{align}
 \delta_t&=r_t+\gamma(1-z_t)V_\psi(o_{t+1}^{\mathrm{final}})-V_\psi(o_t),\\
 \widehat A_t&=\delta_t+\gamma\lambda(1-z_t)(1-\tau_t)\widehat A_{t+1},
 \qquad\lambda=0{,}95.
 \label{eq:gae}
\end{align}
O timeout mantém o bootstrap, mas interrompe a propagação da vantagem através do reset. Se o objetivo fosse vencer uma partida genuinamente finita, tempo restante deveria fazer parte da formulação \cite{pardo2018}. O score no limite de tempo e o objetivo descontado de gols continuam distintos.


\paragraph{A diferença entre terminar o jogo e parar de filmar.}
Um gol encerra o episódio definido pelo treino. Um timeout artificial apenas interrompe a coleta de uma continuação conceitual. Zerar ambos equivale a afirmar que o futuro após qualquer corte vale zero. Para um corte artificial, o erro imediato no alvo TD, se todo o resto for igual, é
\begin{equation}
 \delta_t^{\rm zero}-\delta_t^{\rm bootstrap}=-\gamma V_\psi(o_{t+1}^{\rm final}).
\end{equation}
Se o valor previsto for positivo, zerar o futuro reduz o alvo; se for negativo, aumenta. Por exemplo, $r_t=0$, $V(o_t)=5$, $V(o_{t+1})=5$ e $\gamma=0{,}99$ produzem resíduo $-0{,}05$ com bootstrap, contra $-5$ sem ele. É uma ilustração, não uma medição do bot. Usar a observação após o autoreset seria outro erro: avaliaria uma partida nova como se fosse a continuação da anterior.

A implementação atual distingue essas fronteiras. Ainda assim, prever um futuro errado continua possível. A correção remove uma inconsistência para a semântica escolhida; não torna perfeito o valor, nem transforma o score finito do benchmark no mesmo objetivo do treino.

\section{Atualização PPO e mudanças frente à \texttt{main}}
Com $\rho_t(\theta)=\pi_\theta(a_t\mid o_t)/\pi_{\theta_{\mathrm{old}}}(a_t\mid o_t)$, a implementação segue o objetivo clipped de PPO \cite{schulman2017}:
\begin{align}
 L_{\mathrm{clip}}(\theta)&=\widehat\E_t\!\left[\min\left(
 \rho_t\widehat A_t,\operatorname{clip}(\rho_t,1-\epsilon,1+\epsilon)\widehat A_t\right)\right],\\
 L_{\mathrm{ator}}&=-L_{\mathrm{clip}}-\beta_k\widehat\E_t\mathcal H[\pi_\theta(\cdot\mid o_t)],\\
 L_{\mathrm{crit}}&=\frac{c_V}{2}\widehat\E_t\left[(V_\psi(o_t)-\widehat R_t)^2\right],
 \qquad\epsilon=0{,}2,\quad c_V=0{,}5.
\end{align}
Os parâmetros $\theta$ e $\psi$ são disjuntos. A separação elimina o gradiente direto do crítico nas camadas do ator; $\widehat A_t$ ainda depende de $V_\psi$. Cada grupo é clipado em norma $0{,}5$; a norma conjunta pode atingir $\sqrt2\,0{,}5$.

Nas primeiras $20$ iterações, apenas o crítico é atualizado. O ator carregado permanece invariável e seus gradientes ficam ausentes, evitando passos Adam artificiais. A guarda usa
\begin{equation}
 \widehat D_{\mathrm{KL}}=\widehat\E_t[(\rho_t-1)-\log\rho_t],
 \qquad \widehat D_{\mathrm{KL}}>0{,}03\ \Longrightarrow\ \text{interrupção da atualização restante}.
\end{equation}
É um critério empírico, não uma garantia formal de região de confiança.


\subsection{Por que separar ator e crítico pode estabilizar o warm start}
No tronco compartilhado antigo, sejam $\omega$ seus parâmetros e $g_\pi$, $g_V$ os gradientes de suas perdas. Simplificando para um passo de descida comum,
\begin{equation}
 \Delta\omega=-\alpha(g_\pi+c_Vg_V),\qquad
 \Delta L_\pi\approx-\alpha\left(\norm{g_\pi}^2+c_Vg_\pi^\top g_V\right).
\end{equation}
Se o produto interno for suficientemente negativo, a componente de valor pode contrariar a direção que reduziria a perda da política. A expressão explica uma possibilidade geométrica; o Adam real não é essa atualização simples, e um cosseno negativo isolado não demonstra a causa de uma derrota.

O clipping global também acopla as escalas. Para gradiente concatenado $g$ e limite $c$,
\begin{equation}
 g_{\rm clip}=g\min\left(1,\frac{c}{\norm{g}+\varepsilon}\right).
\end{equation}
Um gradiente grande do crítico pode reduzir junto a componente do ator. Parâmetros e clipping independentes retiram esses dois acoplamentos diretos. Não tornam o ator independente da qualidade do crítico: a vantagem continua sendo sua orientação.

Ao trocar a recompensa, o significado de ``valor futuro'' muda. Preservar um ator que já sabe jogar e permitir vinte iterações de ajuste só do crítico evita mover esse ator durante a primeira recalibração. É semelhante a atualizar a régua de avaliação antes de orientar o jogador com ela. A imobilidade do ator nessa fase foi testada; a superioridade universal dessa duração não foi demonstrada.

\subsection{Por que PPO não promete melhora a cada checkpoint}
O clipping limita o incentivo para mudar muito as probabilidades nos dados coletados; não impõe uma barreira absoluta sobre toda a política. Estados não visitados, reutilização de minibatches, erros de vantagem e adaptação do adversário podem produzir perdas futuras. Mais épocas também não significam automaticamente melhor aprendizado. A guarda de KL e os checkpoints protegem contra alguns riscos operacionais, sem oferecer um teorema de melhoria monotônica do placar.

\subsection{Adversários estáveis e memória de competências}
Contra um adversário $j$, a transição vista pelo aprendiz pode ser escrita como
\begin{equation}
 P_j(s'\mid s,a)=\sum_b P(s'\mid s,a,b)\,\pi_j^-(b\mid O_j(s)).
\end{equation}
Ao alterar o adversário, altera-se essa dinâmica. No modo antigo ele podia mudar entre rollouts, inclusive dentro de uma partida. O crítico precisava acompanhar esse alvo móvel. Na nova mistura, o adversário é sorteado por episódio e permanece fixo nele. O processo aumentado $(s,j)$ tem uma dinâmica estável para uma mistura fixa; como $j$ não é explicitamente observado, o problema do ator ainda pode ser parcialmente observável.

As referências históricas mantêm adversários neurais antigos presentes no treino. O objetivo passa a cobrar, em cada nova amostra de episódio, desempenho na mistura
\begin{equation}
 J_q(\theta)=\sum_j q_jJ_j(\theta).
\end{equation}
Isso torna o esquecimento treinável e mensurável. Entretanto, maximizar uma soma ainda permite trocas entre adversários; não garante melhorar todos. É por isso que os resultados individuais precisam acompanhar a média. Em jogos, uma política pode vencer outra e perder para uma terceira: não existe necessariamente uma escada única de força.

\begin{table}[htbp]
\centering\small
\caption{Diferenças verificadas por inspeção; não constituem ablação causal.}
\begin{tabularx}{\textwidth}{@{}lYY@{}}\toprule
Componente & \texttt{main} inspecionada & Sprint 1 estudada\\\midrule
Recompensa & Padrão da dependência & Equação~\eqref{eq:shaping}\\
Valor/política & Tronco compartilhado & Crítico independente; calibração\\
Adversários & Política atual/snapshot por rollout & Mistura congelada por episódio\\
Timeout & Bootstrap zerado como em término & Equação~\eqref{eq:gae}\\
Desconto & Meia-vida de 5 s & Meia-vida de 12 s\\
Atualizações & Até 4 épocas & Até 6 épocas; 4 minibatches; guarda KL\\
Avaliação & Contra Chase & Dois modos, cenários e referências neurais\\
Retomada & Pesos/snapshots & Adam, schedules, referências, pool e RNGs\\\bottomrule
\end{tabularx}
\end{table}

\section{Resultados observados e unidade de avaliação}
Para vitórias $W$, empates $D$ e derrotas $L$, o score de um caso é
\begin{equation}
 \widehat S=\frac{W+\tfrac12D}{W+D+L}.
\end{equation}
Cada grupo combina adversários, seeds de avaliação e lados com a agregação balanceada do painel. O checkpoint é escolhido por
\begin{equation}
 \widehat J_{\mathrm{sel}}(k)=\min_{m\in\{\mathrm{greedy},\mathrm{sampled}\},\ c\in\{\mathrm{random},\mathrm{kickoff}\}}
 \widehat S_{m,c}(k),\qquad
 k^*=\arg\max_k\widehat J_{\mathrm{sel}}(k).
\end{equation}
O mínimo protege contra um pico em apenas um grupo, mas ainda pode ocultar falhas por adversário.

\begin{table}[htbp]\centering
\caption{Médias descritivas do painel GPU; avaliações sucessivas são correlacionadas.}
\begin{tabular}{@{}lrr@{}}\toprule
Janela de iterações & Registros & Média de $\widehat J_{\mathrm{sel}}$\\\midrule
3.000--4.000 & 21 & 0,86591\\
5.000--6.000 & 22 & 0,86609\\
7.000--8.000 & 21 & 0,87095\\
9.000--10.000 & 21 & 0,88334\\\bottomrule
\end{tabular}
\end{table}
A diferença entre as janelas extremas é aproximadamente $0{,}01743$, ou $1{,}74$ ponto percentual de score. O melhor registro foi $0{,}89644$ na iteração $9\,900$; o endpoint marcou $0{,}88621$. Isso é compatível com progresso oscilatório, não com monotonicidade nem com um teto estabelecido.

O benchmark independente anterior avaliou o checkpoint $2\,550$, não o $9\,900$: foram $3\,324$ casos e $106\,368$ primeiras partidas completas, incluindo repetições determinísticas identificadas. Contra a origem em estados aleatórios, o score foi $0{,}65381$ em greedy e $0{,}66895$ em sampled. A estimativa sampled teve intervalo bootstrap por seed $[0{,}64941;0{,}68799]$. Os resultados desse teste não podem ser atribuídos automaticamente ao modelo GPU.

Repetições idênticas de kickoff greedy contra uma política determinística não geram observações independentes. Seeds de avaliação não substituem seeds de treinamento. Há apenas uma continuação de treinamento neste estudo.


\begin{figure}[htbp]\centering\setlength{\unitlength}{1mm}
\begin{picture}(155,100)
\put(39,10){\line(1,0){100}}
\put(36,88){\makebox(0,0)[r]{\small Origem}}
{\color{black!35}\put(39,90){\rule{50.290mm}{2mm}}}
{\color{blue!65!black}\put(39,87){\rule{75.720mm}{2mm}}}
\put(36,78){\makebox(0,0)[r]{\small Chase}}
{\color{black!35}\put(39,80){\rule{76.570mm}{2mm}}}
{\color{blue!65!black}\put(39,77){\rule{89.160mm}{2mm}}}
\put(36,68){\makebox(0,0)[r]{\small Striker}}
{\color{black!35}\put(39,70){\rule{86.060mm}{2mm}}}
{\color{blue!65!black}\put(39,67){\rule{94.760mm}{2mm}}}
\put(36,58){\makebox(0,0)[r]{\small Keeper}}
{\color{black!35}\put(39,60){\rule{80.820mm}{2mm}}}
{\color{blue!65!black}\put(39,57){\rule{97.120mm}{2mm}}}
\put(36,48){\makebox(0,0)[r]{\small Interceptor}}
{\color{black!35}\put(39,50){\rule{84.910mm}{2mm}}}
{\color{blue!65!black}\put(39,47){\rule{93.710mm}{2mm}}}
\put(36,38){\makebox(0,0)[r]{\small Defensor}}
{\color{black!35}\put(39,40){\rule{88.100mm}{2mm}}}
{\color{blue!65!black}\put(39,37){\rule{96.080mm}{2mm}}}
\put(36,28){\makebox(0,0)[r]{\small Dribbler}}
{\color{black!35}\put(39,30){\rule{86.210mm}{2mm}}}
{\color{blue!65!black}\put(39,27){\rule{95.190mm}{2mm}}}
\put(36,18){\makebox(0,0)[r]{\small Finalizador}}
{\color{black!35}\put(39,20){\rule{74.630mm}{2mm}}}
{\color{blue!65!black}\put(39,17){\rule{95.100mm}{2mm}}}
\put(39,6){\makebox(0,0){\scriptsize 0.0}}
\put(59,6){\makebox(0,0){\scriptsize 0.2}}
\put(79,6){\makebox(0,0){\scriptsize 0.4}}
\put(99,6){\makebox(0,0){\scriptsize 0.6}}
\put(119,6){\makebox(0,0){\scriptsize 0.8}}
\put(139,6){\makebox(0,0){\scriptsize 1.0}}
\put(39,0){\makebox(0,0)[l]{\scriptsize Cinza: política de origem. Azul: checkpoint 2.550.}}
\end{picture}
\caption{Benchmark reservado anterior: score médio com pesos iguais dos quatro grupos modo$\times$cenário. ``Origem na primeira linha é também o adversário fixo. Barras são estimativas agregadas descritivas, sem IC agregado; os ICs por célula estão no relatório auditado. Defensor e Finalizador abreviam CounterDefender e AimedStriker. Não são resultados do checkpoint 9.900.}\label{fig:benchmark}
\end{figure}
\subsection{Como a avaliação antiga podia contar uma história incompleta}
Se um agente obtém score $0{,}90$ contra um perseguidor simples e $0{,}20$ contra uma política neural, avaliar apenas o primeiro esconderia o segundo. Esse exemplo é hipotético; a regressão neural histórica relatada no início é real. Incluir referências e ambos os lados não modifica diretamente o gradiente, mas torna a conclusão e a escolha de checkpoint mais confiáveis.

A política treinada é estocástica. Na sala, a execução greedy escolhe
\begin{equation}
 \mu_\theta(o)=\left(\arg\max_{a_x}\pi_{\theta,x},
 \arg\max_{a_y}\pi_{\theta,y},\arg\max_{a_k}\pi_{\theta,k}\right).
\end{equation}
Essa é outra política. Se duas ações mudam de probabilidades $(0{,}49,0{,}51)$ para $(0{,}51,0{,}49)$, a distribuição muda pouco, mas a ação greedy troca completamente. Em sequências de decisões isso pode mudar a trajetória inteira. O exemplo explica a possibilidade; a inspeção de margens feita no projeto não demonstrou que esse fosse o mecanismo dominante da regressão anterior.

O mínimo dos quatro grupos cobra consistência entre modos e cenários. Porém, nenhum resumo substitui os casos individuais. Um score $0{,}5$ em self-play é especialmente ambíguo: políticas idênticas têm simetria de resultado, sejam fortes ou fracas. No kickoff greedy anterior, tanto o melhor quanto o endpoint empataram todas as partidas simétricas, mas a bola ficou estacionária em $10{,}4\%$ das decisões no primeiro e em $100\%$ no segundo. O placar igual escondia dinâmicas muito diferentes.

\subsection{O que os experimentos permitem atribuir às mudanças}
``Agora deu certo'' tem sustentação operacional: houve ganho em confrontos reservados e o sistema passou a medir regressões antes ocultas. A explicação de mecanismo mais consistente é a combinação de incentivos coerentes, valor recalibrado, diversidade preservada e seleção mais exigente. Não temos, contudo, um experimento longo que isole a contribuição de cada componente.

A ablação curta ajuda a evitar uma narrativa conveniente. Em duas seeds e apenas 40 iterações com 64 ambientes, os endpoints foram:
\begin{table}[htbp]\centering\small
\caption{Ablação exploratória documentada; baseline comum $0{,}69405$. Não é a comparação entre a \texttt{main} inteira e o sistema final.}
\begin{tabular}{@{}lrrr@{}}\toprule
Configuração & Seed 0 & Seed 1 & Média\\\midrule
Corrected, crítico compartilhado & 0,65646 & 0,59226 & 0,62436\\
PBRS, crítico compartilhado & 0,68869 & 0,66084 & 0,67477\\
PBRS, crítico separado & 0,59762 & 0,73750 & 0,66756\\
PBRS, separado e calibração & 0,67560 & 0,72083 & 0,69821\\\bottomrule
\end{tabular}
\end{table}
PBRS superou corrected nos dois pares, mas os orçamentos foram pequenos. Separar o crítico teve efeitos distintos entre seeds. A calibração reservou metade das 40 iterações ao crítico, portanto também mudou o número de atualizações do ator. Os dados não autorizam dizer que cada mudança, isoladamente, sempre melhora o agente.

\begin{figure}[htbp]\centering
\setlength{\unitlength}{1mm}
\begin{picture}(155,76)
\put(18,12){\line(1,0){125}}
\put(18,12){\line(0,1){48}}
\put(16,12.000){\makebox(0,0)[r]{\scriptsize 0.75}}
{\color{black!12}\put(18,12.000){\line(1,0){125}}}
\put(16,24.000){\makebox(0,0)[r]{\scriptsize 0.80}}
{\color{black!12}\put(18,24.000){\line(1,0){125}}}
\put(16,36.000){\makebox(0,0)[r]{\scriptsize 0.85}}
{\color{black!12}\put(18,36.000){\line(1,0){125}}}
\put(16,48.000){\makebox(0,0)[r]{\scriptsize 0.90}}
{\color{black!12}\put(18,48.000){\line(1,0){125}}}
\put(16,60.000){\makebox(0,0)[r]{\scriptsize 0.95}}
{\color{black!12}\put(18,60.000){\line(1,0){125}}}
\put(18.000,8){\makebox(0,0){\scriptsize 3000}}
\put(49.250,8){\makebox(0,0){\scriptsize 4750}}
\put(80.500,8){\makebox(0,0){\scriptsize 6500}}
\put(111.750,8){\makebox(0,0){\scriptsize 8250}}
\put(143.000,8){\makebox(0,0){\scriptsize 10000}}
\put(81,2){\makebox(0,0){\small Iteração de treinamento}}
\put(18,68){\makebox(0,0)[l]{\small Score mínimo dos quatro grupos}}
{\color{blue!65!black}\linethickness{0.32mm}
\qbezier(18.000,39.235)(18.446,39.481)(18.893,39.727)
\qbezier(18.893,39.727)(19.339,37.461)(19.786,35.195)
\qbezier(19.786,35.195)(20.232,37.959)(20.679,40.723)
\qbezier(20.679,40.723)(21.125,40.366)(21.571,40.010)
\qbezier(21.571,40.010)(22.018,40.287)(22.464,40.564)
\qbezier(22.464,40.564)(22.911,40.197)(23.357,39.831)
\qbezier(23.357,39.831)(23.804,39.751)(24.250,39.671)
\qbezier(24.250,39.671)(24.696,40.458)(25.143,41.244)
\qbezier(25.143,41.244)(25.589,40.782)(26.036,40.320)
\qbezier(26.036,40.320)(26.482,40.421)(26.929,40.523)
\qbezier(26.929,40.523)(27.375,38.366)(27.821,36.210)
\qbezier(27.821,36.210)(28.268,39.099)(28.714,41.988)
\qbezier(28.714,41.988)(29.161,40.569)(29.607,39.149)
\qbezier(29.607,39.149)(30.054,39.392)(30.500,39.634)
\qbezier(30.500,39.634)(30.946,40.188)(31.393,40.742)
\qbezier(31.393,40.742)(31.839,40.508)(32.286,40.274)
\qbezier(32.286,40.274)(32.732,40.036)(33.179,39.798)
\qbezier(33.179,39.798)(33.625,39.448)(34.071,39.098)
\qbezier(34.071,39.098)(34.518,40.364)(34.964,41.629)
\qbezier(34.964,41.629)(35.411,41.132)(35.857,40.634)
\qbezier(35.857,40.634)(36.304,40.706)(36.750,40.779)
\qbezier(36.750,40.779)(37.196,41.609)(37.643,42.440)
\qbezier(37.643,42.440)(38.089,42.539)(38.536,42.639)
\qbezier(38.536,42.639)(38.982,42.117)(39.429,41.594)
\qbezier(39.429,41.594)(39.875,41.977)(40.321,42.359)
\qbezier(40.321,42.359)(40.768,39.527)(41.214,36.695)
\qbezier(41.214,36.695)(41.661,38.435)(42.107,40.175)
\qbezier(42.107,40.175)(42.554,41.602)(43.000,43.028)
\qbezier(43.000,43.028)(43.446,41.800)(43.893,40.572)
\qbezier(43.893,40.572)(44.339,41.090)(44.786,41.609)
\qbezier(44.786,41.609)(45.232,41.842)(45.679,42.075)
\qbezier(45.679,42.075)(46.125,40.585)(46.571,39.095)
\qbezier(46.571,39.095)(47.018,40.331)(47.464,41.568)
\qbezier(47.464,41.568)(47.911,39.000)(48.357,36.433)
\qbezier(48.357,36.433)(48.804,38.683)(49.250,40.933)
\qbezier(49.250,40.933)(49.696,41.973)(50.143,43.013)
\qbezier(50.143,43.013)(50.589,42.995)(51.036,42.976)
\qbezier(51.036,42.976)(51.482,43.147)(51.929,43.318)
\qbezier(51.929,43.318)(52.375,44.017)(52.821,44.715)
\qbezier(52.821,44.715)(53.268,43.342)(53.714,41.970)
\qbezier(53.714,41.970)(54.161,42.528)(54.607,43.086)
\qbezier(54.607,43.086)(55.054,42.545)(55.500,42.004)
\qbezier(55.500,42.004)(55.946,42.120)(56.393,42.235)
\qbezier(56.393,42.235)(56.839,43.176)(57.286,44.118)
\qbezier(57.286,44.118)(57.732,44.311)(58.179,44.504)
\qbezier(58.179,44.504)(58.625,31.615)(59.071,18.726)
\qbezier(59.071,18.726)(59.518,27.677)(59.964,36.629)
\qbezier(59.964,36.629)(60.411,27.466)(60.857,18.304)
\qbezier(60.857,18.304)(61.304,30.492)(61.750,42.680)
\qbezier(61.750,42.680)(62.196,42.579)(62.643,42.478)
\qbezier(62.643,42.478)(63.089,42.837)(63.536,43.195)
\qbezier(63.536,43.195)(63.982,43.950)(64.429,44.706)
\qbezier(64.429,44.706)(64.875,43.223)(65.321,41.740)
\qbezier(65.321,41.740)(65.768,38.938)(66.214,36.137)
\qbezier(66.214,36.137)(66.661,36.308)(67.107,36.479)
\qbezier(67.107,36.479)(67.554,38.772)(68.000,41.065)
\qbezier(68.000,41.065)(68.446,41.830)(68.893,42.596)
\qbezier(68.893,42.596)(69.339,42.671)(69.786,42.746)
\qbezier(69.786,42.746)(70.232,43.252)(70.679,43.757)
\qbezier(70.679,43.757)(71.125,43.415)(71.571,43.073)
\qbezier(71.571,43.073)(72.018,43.900)(72.464,44.728)
\qbezier(72.464,44.728)(72.911,40.554)(73.357,36.381)
\qbezier(73.357,36.381)(73.804,39.923)(74.250,43.466)
\qbezier(74.250,43.466)(74.696,43.560)(75.143,43.655)
\qbezier(75.143,43.655)(75.589,43.337)(76.036,43.019)
\qbezier(76.036,43.019)(76.482,43.055)(76.929,43.090)
\qbezier(76.929,43.090)(77.375,43.612)(77.821,44.134)
\qbezier(77.821,44.134)(78.268,43.256)(78.714,42.377)
\qbezier(78.714,42.377)(79.161,39.172)(79.607,35.967)
\qbezier(79.607,35.967)(80.054,39.969)(80.500,43.970)
\qbezier(80.500,43.970)(80.946,40.150)(81.393,36.329)
\qbezier(81.393,36.329)(81.839,39.943)(82.286,43.557)
\qbezier(82.286,43.557)(82.732,44.531)(83.179,45.506)
\qbezier(83.179,45.506)(83.625,44.840)(84.071,44.174)
\qbezier(84.071,44.174)(84.518,44.239)(84.964,44.304)
\qbezier(84.964,44.304)(85.411,44.322)(85.857,44.339)
\qbezier(85.857,44.339)(86.304,44.157)(86.750,43.975)
\qbezier(86.750,43.975)(87.196,43.931)(87.643,43.887)
\qbezier(87.643,43.887)(88.089,44.710)(88.536,45.533)
\qbezier(88.536,45.533)(88.982,45.680)(89.429,45.828)
\qbezier(89.429,45.828)(89.875,43.999)(90.321,42.171)
\qbezier(90.321,42.171)(90.768,42.729)(91.214,43.287)
\qbezier(91.214,43.287)(91.661,43.324)(92.107,43.362)
\qbezier(92.107,43.362)(92.554,39.707)(93.000,36.053)
\qbezier(93.000,36.053)(93.446,40.522)(93.893,44.992)
\qbezier(93.893,44.992)(94.339,43.453)(94.786,41.914)
\qbezier(94.786,41.914)(95.232,38.965)(95.679,36.016)
\qbezier(95.679,36.016)(96.125,36.215)(96.571,36.415)
\qbezier(96.571,36.415)(97.018,40.603)(97.464,44.792)
\qbezier(97.464,44.792)(97.911,44.019)(98.357,43.247)
\qbezier(98.357,43.247)(98.804,43.023)(99.250,42.799)
\qbezier(99.250,42.799)(99.696,44.120)(100.143,45.441)
\qbezier(100.143,45.441)(100.589,45.452)(101.036,45.462)
\qbezier(101.036,45.462)(101.482,40.893)(101.929,36.324)
\qbezier(101.929,36.324)(102.375,40.693)(102.821,45.062)
\qbezier(102.821,45.062)(103.268,44.786)(103.714,44.509)
\qbezier(103.714,44.509)(104.161,40.172)(104.607,35.835)
\qbezier(104.607,35.835)(105.054,27.348)(105.500,18.861)
\qbezier(105.500,18.861)(105.946,31.639)(106.393,44.416)
\qbezier(106.393,44.416)(106.839,44.605)(107.286,44.793)
\qbezier(107.286,44.793)(107.732,44.641)(108.179,44.490)
\qbezier(108.179,44.490)(108.625,44.614)(109.071,44.739)
\qbezier(109.071,44.739)(109.518,44.428)(109.964,44.117)
\qbezier(109.964,44.117)(110.411,44.393)(110.857,44.669)
\qbezier(110.857,44.669)(111.304,31.788)(111.750,18.908)
\qbezier(111.750,18.908)(112.196,31.381)(112.643,43.854)
\qbezier(112.643,43.854)(113.089,44.383)(113.536,44.912)
\qbezier(113.536,44.912)(113.982,44.704)(114.429,44.496)
\qbezier(114.429,44.496)(114.875,44.887)(115.321,45.278)
\qbezier(115.321,45.278)(115.768,44.071)(116.214,42.863)
\qbezier(116.214,42.863)(116.661,44.543)(117.107,46.223)
\qbezier(117.107,46.223)(117.554,45.580)(118.000,44.936)
\qbezier(118.000,44.936)(118.446,45.307)(118.893,45.677)
\qbezier(118.893,45.677)(119.339,45.120)(119.786,44.564)
\qbezier(119.786,44.564)(120.232,45.140)(120.679,45.716)
\qbezier(120.679,45.716)(121.125,44.848)(121.571,43.980)
\qbezier(121.571,43.980)(122.018,44.962)(122.464,45.943)
\qbezier(122.464,45.943)(122.911,41.171)(123.357,36.399)
\qbezier(123.357,36.399)(123.804,40.674)(124.250,44.949)
\qbezier(124.250,44.949)(124.696,40.288)(125.143,35.628)
\qbezier(125.143,35.628)(125.589,40.545)(126.036,45.462)
\qbezier(126.036,45.462)(126.482,44.659)(126.929,43.855)
\qbezier(126.929,43.855)(127.375,44.150)(127.821,44.445)
\qbezier(127.821,44.445)(128.268,44.879)(128.714,45.312)
\qbezier(128.714,45.312)(129.161,40.694)(129.607,36.076)
\qbezier(129.607,36.076)(130.054,40.933)(130.500,45.789)
\qbezier(130.500,45.789)(130.946,45.133)(131.393,44.478)
\qbezier(131.393,44.478)(131.839,44.411)(132.286,44.345)
\qbezier(132.286,44.345)(132.732,45.581)(133.179,46.817)
\qbezier(133.179,46.817)(133.625,45.866)(134.071,44.915)
\qbezier(134.071,44.915)(134.518,44.416)(134.964,43.917)
\qbezier(134.964,43.917)(135.411,44.563)(135.857,45.210)
\qbezier(135.857,45.210)(136.304,46.051)(136.750,46.893)
\qbezier(136.750,46.893)(137.196,46.152)(137.643,45.412)
\qbezier(137.643,45.412)(138.089,45.824)(138.536,46.236)
\qbezier(138.536,46.236)(138.982,41.407)(139.429,36.577)
\qbezier(139.429,36.577)(139.875,41.120)(140.321,45.663)
\qbezier(140.321,45.663)(140.768,46.405)(141.214,47.146)
\qbezier(141.214,47.146)(141.661,46.153)(142.107,45.161)
\qbezier(142.107,45.161)(142.554,44.925)(143.000,44.690)
}\end{picture}
\caption{Curva real do painel GPU, sem suavização. Em iterações repetidas após retomada, usa-se o último registro. O eixo vertical começa em 0,75 para mostrar oscilações; não representa taxa de vitória. Painel reutilizado para seleção: os pontos não são replicações independentes e a curva não é validação reservada do checkpoint 9.900.}\label{fig:gpu}
\end{figure}

\section{Diagnóstico da desaceleração}
\subsection{Schedules e exploração: candidatos, não causas demonstradas}
As médias de duas janelas de mil iterações são apresentadas na Tabela~\ref{tab:diag}.
\begin{table}[htbp]\centering\small
\caption{Telemetria da continuação GPU.}\label{tab:diag}
\begin{tabular}{@{}lrr@{}}\toprule
Grandeza & 3.001--4.000 & 9.001--10.000\\\midrule
Taxa de aprendizado & $2{,}30003\,10^{-4}$ & $1{,}09991\,10^{-4}$\\
Entropia (nats) & 1,35034 & 1,20892\\
KL aproximada & 0,003180 & 0,001898\\
Fração de clipping & 0,03426 & 0,01686\\
Variância explicada & 0,84374 & 0,87054\\
Fração de bola estacionária & 0,21200 & 0,21301\\
Recompensa média/transição & 0,28196 & 0,33881\\
Interrupções por guarda KL & 0/1.000 & 0/1.000\\\bottomrule
\end{tabular}
\end{table}
O schedule é
\begin{equation}
 \alpha_k=3\,10^{-4}(1-f_k)+10^{-4}f_k,\quad
 \beta_k=0{,}02(1-f_k)+0{,}01f_k,\quad
 f_k=\frac{k-1}{K-1},\quad K=10\,000.
\end{equation}
Passos menores podem reduzir progresso e também evitar instabilidade. A guarda de KL não foi acionada nessas janelas; os dados não sustentam responsabilizá-la pelo freio. Para as três cabeças,
\begin{equation}
 \mathcal H_{\max}=\log3+\log3+\log2=\log18\approx2{,}8904.
\end{equation}
A entropia diminuiu, mas não demonstra colapso total. Variância explicada alta tampouco certifica precisão em situações raras:
\begin{equation}
 \operatorname{EV}=1-\frac{\operatorname{Var}(\widehat R-V_\psi)}{\operatorname{Var}(\widehat R)}.
\end{equation}

\subsection{Currículo, observação e crédito temporal}
A mistura por novo episódio permanece em $65\%$ de heurísticas e $35\%$ de referências neurais fixas. Um conjunto finito pode deixar de pressionar novas respostas; isso é hipótese a testar com oponentes recentes, mantendo âncoras contra esquecimento. Mudança de adversário pode produzir mais aprendizado que ajuste de recompensa, mas ainda não houve esse contraste controlado.

A política não tem memória explícita. Isso pode limitar decisões dependentes de informações internas não observadas; não prova que a dimensão da rede esteja insuficiente. No GAE, o fator de propagação é $\gamma\lambda\approx0{,}94271$, com escala geométrica indicativa
\begin{equation}
 \frac{1}{1-\gamma\lambda}\approx17{,}45\ \text{decisões}\approx2{,}33\ \mathrm{s}.
\end{equation}
Essa escala não é um horizonte rígido: o bootstrap do crítico pode representar consequências mais longas. Ainda assim, jogadas encadeadas exigem estimativas de valor adequadas para atribuir mérito a ações anteriores.

\subsection{Eficiência computacional}
Na janela final GPU, coleta e atualização levaram aproximadamente $1{,}102$ s e $0{,}336$ s. A coleta inclui simulação, inferência e transferências de dados. Para fração acelerável $f\approx0{,}234$ e aceleração $a$ apenas da atualização, a decomposição idealizada fornece
\begin{equation}
 S(a)=\frac{1}{(1-f)+f/a},\qquad S(\infty)\approx1{,}31.
\end{equation}
Esse limite ilustrativo ignora avaliações e outros custos; não prevê uma máquina específica. Mostra que acelerar somente a atualização já não acelera todo o treinamento proporcionalmente. Tempos de máquinas distintas não isolam causalmente o efeito da GPU.

\section{Protocolo mínimo para testar novos incentivos}
A variante de contenção do mapa alterou colisões nas fronteiras, preservando parâmetros e geometria. É necessário estabelecer uma referência nova nesse mapa. Os $2\,160$ testes físicos sem fuga não são um benchmark de habilidade.

Cada braço deve partir do mesmo ator $\theta^*$, usar orçamento e mapa iguais e alterar um único fator. Em warm starts, o controle também deve reiniciar Adam e receber a mesma calibração do crítico; do contrário, recompensa e reinicialização ficam confundidas. Uma comparação válida é
\begin{equation}
 \Delta_z=S_{\mathrm{variante},z}-S_{\mathrm{controle},z},\qquad
 \widehat\Delta=\frac{1}{Z}\sum_{z=1}^{Z}\Delta_z.
\end{equation}
Reamostrar blocos de seeds preserva o pareamento; intervalos são descritivos e não eliminam a necessidade de seeds de treinamento independentes. Recompensas acumuladas de tabelas diferentes não devem ser diretamente usadas para eleger uma vencedora. Adversários e seeds reservados devem permanecer fora da seleção repetida de checkpoints.


\section{Como levar o raciocínio a Pokémon, robótica e outros jogos}
Os exemplos desta seção são propostas de desenho, não resultados obtidos pelo nosso modelo. A transferência mais importante não é copiar os pesos $0{,}15$ e $0{,}08$. É repetir o raciocínio: definir sucesso, separar sucesso de pistas, entender o que o agente observa e testar contra desafios que não serviram para escolher o modelo.

\subsection{Pokémon: ganhar uma batalha não é apenas causar dano}
Antes de escrever uma recompensa, é preciso escolher a tarefa. Um agente de batalha, um agente que percorre o mundo e um agente que tenta completar uma coleção resolvem problemas diferentes. Para batalhas, uma referência simples poderia ser $+1$ por vitória e $-1$ por derrota. Para exploração, o objetivo pode estar ligado a marcos de progresso. Não existe uma única tabela ideal para tudo que chamamos de ``jogar Pokémon''.

Considere o agente de batalha. Premiar dano causado parece sensato, mas uma ação que causa dano agora pode desperdiçar um recurso ou expor uma peça essencial. Uma troca sem dano imediato pode ser a decisão que permite ganhar depois. Assim como perseguir a bola não é igual a jogar bem HaxBall, causar dano não é igual a maximizar chance de vitória.

Como exemplo inicial de potencial, poderíamos usar uma nota limitada de recursos do time:
\begin{equation}
 \Phi(s)=c\left[\frac{1}{n}\sum_{i=1}^{n}h_i^{\rm nosso}
 -\frac{1}{m}\sum_{j=1}^{m}h_j^{\rm adversario}\right],
 \qquad 0\leq h\leq1.
\end{equation}
Aqui $h$ representa a fração de vida de cada membro de uma composição fixa, incluindo derrotados com valor zero. É apenas uma heurística: pesos iguais ignoram funções estratégicas, composição e informação desconhecida. Mesmo que a fórmula facilite o começo, não é uma avaliação confiável de todas as posições. Em observação parcial, nem todos esses valores estarão disponíveis ao agente; usar informação privilegiada no treino precisa ser uma decisão explícita, especialmente para não vazar dados ao ator na avaliação.

Dar um bônus a cada turno por ter muita vida e aplicar a diferença descontada desse potencial são tratamentos distintos. O primeiro pode valorizar prolongar a batalha ou recuperar vida sem progredir; o segundo exige fechar a conta. Todos os terminais relevantes precisam receber o tratamento consistente, inclusive derrotas. Usar potencial zero no fim apenas das vitórias deixaria uma fronteira diferente nas derrotas.

A informação também muda o problema. No nosso HaxBall, posição e velocidade são importantes. Em batalhas com recursos ocultos, o histórico pode informar hipóteses sobre o adversário. Uma rede sem memória pode receber duas observações aparentemente iguais que exigem decisões diferentes devido ao passado. Aumentar o bônus de dano não resolve essa falta de informação. Uma memória ou um resumo explícito do histórico seria uma hipótese arquitetural a comparar.

O benchmark deveria separar adversários, composições e estilos conhecidos de combinações reservadas. Um agente pode decorar uma população estreita. A pergunta útil é se a vitória se mantém em novas situações, e não apenas se a recompensa de dano continua subindo.

\subsection{Pokémon de exploração: o exemplo da porta trancada}
Imagine um objetivo atrás de uma porta que exige buscar uma chave longe dali. A distância em linha reta ao objetivo piora quando o agente vai pegar a chave. Um bônus imediato por diminuir a distância pode desincentivar exatamente a preparação necessária.

Um potencial de distância descontado pode preservar o objetivo sob suas condições, mas ainda oferecer orientação local ruim. Uma nota que represente também possuir a chave e acessar novas regiões seria uma hipótese melhor informada. Isso mostra duas distinções: uma pista pode ser formalmente compatível com o objetivo e, mesmo assim, pouco útil; e uma observação que não representa a chave pode impedir o agente de distinguir situações essenciais.

Também é preciso decidir se obter a chave é uma pista para a missão ou parte do objetivo em si. Se o usuário deseja explorar e colecionar, recompensar descobertas pode definir legitimamente a tarefa. Nesse caso, não há motivo para alegar que se preservou um objetivo antigo de apenas chegar ao destino: decidiu-se treinar outra preferência.

\subsection{Robótica: aproximar a mão não é o mesmo que concluir a tarefa}
Considere, como exemplo conceitual em simulação, um braço que precisa pegar um objeto e colocá-lo numa região. Uma recompensa por proximidade pode ensinar aproximação sem preensão; uma recompensa por elevar o objeto pode não ensinar depositá-lo corretamente. As etapas não são intercambiáveis com o sucesso completo.

Um potencial simples seria $\Phi(s)=-c\,d(s)$, com distância limitada ao alvo. Mas que alvo? No começo, a mão precisa alcançar o objeto; depois, o objeto precisa alcançar a região de entrega. Essa escolha exige representar a fase da tarefa no estado e testar como as transições de fase entram no potencial. Trocar manualmente a regra no meio da execução sem registrar o contexto pode produzir sinais contraditórios.

Se o objetivo inclui suavidade ou baixo consumo, uma formulação possível é
\begin{equation}
 r_{\rm tarefa}=r_{\rm sucesso}-\lambda_E\|u_t\|^2
 -\lambda_S\|u_t-u_{t-1}\|^2,
\end{equation}
onde $u_t$ é o comando. Esses custos mudam o objetivo: aceitam trocar sucesso, esforço e suavidade conforme seus pesos. Não são automaticamente PBRS. Certos limites operacionais podem precisar ser restrições da tarefa ou do controlador; preservar uma política ótima de recompensa não estabelece uma garantia física de execução.

Há ainda uma particularidade temporal. Se passarmos de comandos a cada $0{,}1$ s para comandos a cada $0{,}01$ s, copiar o mesmo $\gamma$ por decisão muda o quanto o futuro pesa em segundos. Uma parametrização por meia-vida física evita essa troca silenciosa:
\begin{equation}
 \gamma(\Delta t)=2^{-\Delta t/T_{1/2}}.
\end{equation}
A ideia já aparece no nosso treino de HaxBall. Em jogos por turnos, pode fazer mais sentido pensar em turnos; em controle físico, em tempo real. Se a duração das ações varia, o desconto precisa refletir a duração escolhida na formulação.

Uma alternativa de pesquisa é o Hindsight Experience Replay, que reaproveita experiências considerando metas efetivamente alcançadas, estudado em tarefas com objetivos \cite{andrychowicz2017}. Ele não é a mesma técnica que PBRS. A versão clássica foi desenvolvida com aprendizado off-policy; não se deve anexar replay de metas ao nosso PPO on-policy e presumir compatibilidade sem mudanças no algoritmo.

\subsection{Outros jogos competitivos: o adversário também faz parte do problema}
Em corrida, velocidade instantânea não substitui completar a pista; em estratégia, coletar recursos não substitui vencer; em um jogo de luta, causar dano não substitui administrar posição e risco. São analogias de desenho, não diagnósticos de implementações específicas. Em cada caso, um indicador intermediário pode ser útil para explicar progresso sem necessariamente virar recompensa.

Contra um adversário congelado, podemos estudar o ambiente induzido por ele. Contra um oponente que aprende junto, esse ambiente muda. O telescopamento do potencial continua sendo uma conta de trajetória quando suas condições valem; não torna estacionária a população de oponentes nem garante convergência de self-play.

A identidade também ajuda a pensar no caso multiagente: se, para cada jogador, o retorno de qualquer combinação de estratégias muda só por uma constante independente das estratégias, suas preferências entre respostas permanecem. Esse é um argumento por identidade de retorno, não uma alegação de que nosso PPO encontre um equilíbrio. É preciso fechar as fronteiras para todos os jogadores e manter clara a informação usada por cada um.

A lição das âncoras do HaxBall se transporta bem como hipótese: continuar enfrentando estilos antigos ajuda a medir e treinar retenção. Depois, desafios novos podem ampliar a pressão de aprendizado. Atualizar todos os adversários, a recompensa e a observação ao mesmo tempo dificulta saber por que o resultado mudou.

\begin{table}[htbp]\centering\small
\caption{O mesmo raciocínio em tarefas diferentes. Exemplos prospectivos; não são configurações já validadas.}
\begin{tabularx}{\textwidth}{@{}lYYY@{}}\toprule
Caso & Sucesso de referência & Pista possível & O que a pista não prova\\\midrule
HaxBall & Retorno de gols & Posição atrás da bola & Que defender ou recuar seja ruim\\
Batalha & Vencer & Balanço de recursos & Que dano imediato seja a melhor jogada\\
Exploração & Cumprir missão & Acesso e distância & Que o caminho mais curto esteja liberado\\
Manipulação & Pegar e depositar & Distância por fase & Que aproximação signifique preensão\\
Corrida & Completar a tarefa & Avanço na pista & Que velocidade alta signifique controle\\\bottomrule
\end{tabularx}
\end{table}

\subsection{Um roteiro de perguntas para novos projetos}
Primeiro, escreva o que contaria como sucesso em uma frase concreta. Depois, pergunte o que o agente consegue observar e o que precisa lembrar. Só então escolha uma pista intermediária e decida se ela deve mudar o objetivo ou apenas orientar o aprendizado. Se a justificativa for PBRS, confira desconto, estado inicial, terminais e cortes artificiais.

Antes de investir num treino longo, construa exemplos pequenos: ficar parado, avançar e recuar, repetir um ciclo, falhar e alcançar a meta. A soma auxiliar deve bater com a identidade esperada. Em seguida, verifique que a avaliação mede a tarefa original, com situações reservadas. Se mudar o objetivo, essa mudança deve aparecer explicitamente na documentação e na comparação.

Esse roteiro não fornece os melhores pesos. Ele evita confundir mais recompensa com mais competência. No HaxBall, essa distinção foi essencial para deixar de interpretar qualquer pico como uma solução.

\section{Perspectiva posterior: cooperação 2v2}
\subsection{Contratos que precisam ser corrigidos}
No builder instalado,
\begin{equation}
 d(P)=12+4(P-1),\qquad d(2)=16,\quad d(4)=24.
\end{equation}
O teste de reset retornou
\begin{equation}
 (p_{R_1},p_{R_2},p_{B_1},p_{B_2})=
 ((-200,0),(200,0),(0,0),(0,0)).
\end{equation}
Há um aliado no lado oposto e dois jogadores sobrepostos. A ordem dos demais jogadores também muda o significado dos slots entre times. Antes de treinar, precisamos de spawns válidos, organização consistente de aliados/adversários, testes de permutação e simetria, ações para quatro jogadores e exportação compatível. O preset atual não está validado para 2v2.

\subsection{Transferência do ator com semântica explícita}
Se $W\in\R^{256\times16}$ é a primeira camada e $\widetilde o\in\R^{24}$ a nova observação, uma transferência inicial possível é
\begin{equation}
 \widetilde W=WP,\qquad P\in\R^{16\times24},\qquad
 \widetilde W\widetilde o+b=W(P\widetilde o)+b.
\end{equation}
$P$ deve selecionar grandezas equivalentes, incluindo um adversário de referência. Colocar o aliado no slot antigo do adversário e copiar pesos por posição seria semanticamente incorreto. Novas colunas podem iniciar em zero, com posterior ajuste de todas as camadas. O crítico e o otimizador exigem tratamento explícito; o loader atual rejeita corretamente dimensões incompatíveis. Devem ser comparados transferência e treino do zero.

\subsection{Objetivo coletivo e apoio sem bola}
Uma proposta é compartilhar o ator entre companheiros e usar crítico com estado global durante o treino, mantendo execução por observação individual. Essa separação é estudada em MAPPO \cite{yu2022}. Não se presume sua superioridade neste ambiente.

Para adversários congelados, a política conjunta pode ser fatorada como
\begin{equation}
 \pi_\theta^{\mathcal T}(a^1,a^2\mid o^1,o^2)
 =\pi_\theta(a^1\mid o^1)\pi_\theta(a^2\mid o^2),
 \qquad V_\psi=V_\psi(s,i).
\end{equation}
A recompensa de gol é coletiva. Repetir para ambos o incentivo de ficar atrás da bola pode orientar a exploração para aglomeração; isso não torna cooperação teoricamente impossível. Como hipótese de shaping conjunto, sejam $q_a(s)$ um alvo de pressão e $q_s(s)$ um alvo de apoio lateral/cobertura:
\begin{equation}
 \Phi_{\mathcal T}(s)=-\frac{1}{\ell}\min_{\sigma\in\mathfrak S_2}
 \left[w_a\min(D,\norm{p_{\sigma(1)}-q_a})+
 w_s\min(D,\norm{p_{\sigma(2)}-q_s})\right].
\end{equation}
A minimização permite trocar papéis sem fixar um jogador como atacante permanente. Os alvos e pesos ainda precisam de especificação e validação. Com potencial limitado, desconto consistente e terminal zero,
\begin{equation}
 r_t^{\mathcal T}=25(g_t^+-g_t^-)+
 \eta[\gamma\bar\Phi_{\mathcal T}(s_{t+1})-\Phi_{\mathcal T}(s_t)]
\end{equation}
herda a identidade telescópica para o retorno coletivo. Isso é uma proposta, não uma implementação ou resultado experimental.

\subsection{Currículo de passes e critérios de sucesso}
Um passe completo deve ser instrumentado como toque real de $A$ seguido de controle por $B$, sem toque adversário intermediário. Com limiares pré-especificados,
\begin{equation}
 I_{\mathrm{passe}}=\ind\{\ell_t=A,\ \ell_{t'}=B,\
 0<t'-t\leq\Delta_p,\ \norm{b_{t'}-b_t}\geq d_p,\ C_B(t')=1\},
\end{equation}
onde $\ell$ representa eventos de toque e $C_B$ controle sustentado. O estado disponível não expõe diretamente esses eventos; proximidade geométrica não basta. A métrica tampouco identifica intenção. Inicialmente, medir passes é preferível a premiar qualquer troca, que pode gerar circulação sem utilidade.

O estudo deve avançar por: (i) passe para receptor livre em posições variadas; (ii) pressão em 2v1 com alternativas de condução/finalização; (iii) passe seguido de deslocamento para nova linha; (iv) 2v2 com duplas e parceiros diversos. Medir recepção, perdas, cobertura, gols após passe e casos em que passar é uma decisão ruim. Validar também troca de lados e IDs. Nenhuma dessas fases foi treinada neste relatório.

\section{Síntese: por que o conjunto passou a funcionar melhor}
A evidência é restrita a uma continuação, a um painel reutilizado e a benchmarks específicos. A comparação com a \code{main} não é causal; ganhos do checkpoint anterior não certificam o atual. Os dados não demonstram teto definitivo, planejamento de jogadas ou força humana consolidada.

A interpretação mais defensável é que o sistema anterior aprendia competências parciais com orientação e medição frágeis. A mudança estrutural alinhou a soma da recompensa ao objetivo descontado de gols, retirou interferências diretas do crítico no ator, manteve adversários de referência e passou a selecionar modelos sob condições mais amplas. Há razões matemáticas para essas decisões e evidência empírica de progresso do conjunto; falta identificar causalmente quanto cada decisão contribuiu.

Para o leitor que busca uma resposta direta: não foi simplesmente ``colocar mais pontos no gol'' nem ``treinar por mais tempo''. Foi reorganizar como o agente recebe orientação, estima consequências, pratica e é examinado. O resultado melhorou nos testes registrados, mas permanece vulnerável a situações e adversários não cobertos.

Os próximos testes devem distinguir diversidade de desafios, schedules, exploração e observação, sem atribuir toda desaceleração à recompensa. Para 2v2, o avanço imediato é corrigir os contratos e demonstrar recepção e apoio em cenários controlados. Não houve alteração de pesos nem novo treinamento durante este estudo.

\appendix
\section{Rastreabilidade e leitura dos indicadores}
Os dados numéricos foram extraídos de \code{experiments/sprint1_structural/gpu_hourly/training.jsonl} e \code{evaluations.jsonl}. Janelas de treino contêm $1\,000$ registros cada; as janelas de avaliação têm $21$, $22$, $21$ e $21$ registros, exatamente como armazenados. Não são amostras independentes de treinamento.

A extração reprodutível está em \code{articles/evolucao_limites_e_2v2_evidence.json}; o benchmark anterior, em \code{experiments/sprint1_structural/benchmark_3000}; a correção física, em \code{experiments/map_containment}. As aplicações fora de HaxBall e os exemplos de ciclos são construções didáticas, não medições. Este documento é autônomo: não depende de bibliografia ou figuras externas para compilar.

\begin{thebibliography}{9}
\bibitem{ng1999} A. Y. Ng, D. Harada e S. Russell.
\emph{Policy Invariance Under Reward Transformations: Theory and Application to Reward Shaping}. ICML, 1999.
\url{https://people.eecs.berkeley.edu/~pabbeel/cs287-fa09/readings/NgHaradaRussell-shaping-ICML1999.pdf}.
\bibitem{schulman2017} J. Schulman, F. Wolski, P. Dhariwal, A. Radford e O. Klimov.
\emph{Proximal Policy Optimization Algorithms}. 2017.
\url{https://arxiv.org/abs/1707.06347}.
\bibitem{pardo2018} F. Pardo, A. Tavakoli, V. Levdik e P. Kormushev.
\emph{Time Limits in Reinforcement Learning}. ICML, PMLR 80, pp.~4045--4054, 2018.
\url{https://proceedings.mlr.press/v80/pardo18a.html}.
\bibitem{yu2022} C. Yu et al.
\emph{The Surprising Effectiveness of PPO in Cooperative, Multi-Agent Games}. 2022.
\url{https://arxiv.org/abs/2103.01955}.
\bibitem{andrychowicz2017} M. Andrychowicz et al.
\emph{Hindsight Experience Replay}. NeurIPS, 2017.
\url{https://arxiv.org/abs/1707.01495}.
\end{thebibliography}
\end{document}
